\chapter{Introduction}
\label{ch:intro_pyclical}

\section{What is PyClical?}
PyClical is a Python extension module that is designed to make it easy for you to use Python to perform computations in Clifford algebras, and to implement algorithms and programs that use these computations. The syntax and semantics of the main \texttt{index\_set} and \texttt{clifford} classes in PyClical are designed to be as similar as possible to that of the \glucat template classes \texttt{index\_set<>} and \texttt{matrix\_multi<double>}, so that you should be able to use PyClical to teach yourself how to use calculations in Clifford algebras, and then transfer that knowledge to your ability to use \glucat template classes in C++.

\section{Prerequisites}
To make the best use of the \glucat template library and the PyClical extension module, you will need to become familiar with Geometric Algebra and Clifford algebras in general.
Recommended resources include:
\begin{itemize}
    \item Wikipedia: \href{http://en.wikipedia.org/wiki/Geometric_algebra}{Geometric algebra}
    \item Alan Macdonald, \href{http://faculty.luther.edu/~macdonal/GA&GC.pdf}{A Survey of Geometric Algebra and Geometric Calculus}
\end{itemize}
